% Complete reviewed chat addendum, reconstructed after the outage.
% Not a new numbered theorem and not an all-depth construction.
% Owner and cohomology_audit checked the catalog and full arithmetic proof.
% Literature independently checked the displayed added-field/root gluing,
% including the explicit lower-root marking correction.
% No new numerical computation is claimed.

\subsubsection*{The common twenty-five-member class-four catalog}
Let $S_0=\{2,5,83,173,\infty\}$ and let $t_0=83/256$.
For every E.215 tail, all fresh valuation exponents are divisible by five.
Its integral matched class-two section reduces, after one global central
conjugation, to the reference class-two section modulo five. Indeed it is
unramified outside $S_0$ in that quotient, and
\[
                 H^1(G_{\mathbb Q,S_0},\mathbb F_5(2))=0.
\]
The fixed-$\{5,\infty\}$ group is zero, and every added prime is
primitive modulo five, so its possible tame residue weight-one invariant
space is zero. This is a statement about the actual finite section and
its cocycles, not only about the two abelian cohomology classes.

Let $M_0=\Gamma_3/\Gamma_5$ with the actual action twisted by the
reference class-two section, and let $\overline M=M_0/5$. This is the same
finite module for every E.215 tail. It fits into
\[
 0\longrightarrow\mathbb Z_5(4)^3\longrightarrow M_0
 \longrightarrow\mathbb Z_5(3)^2\longrightarrow0.
\]
The fixed-$S_0$ second cohomology groups in weights three and four vanish:
the fixed-five groups vanish by E.209, and the extra tame residue weights
two and three have no invariants at primitive primes. Consequently
$H^2(G_{\mathbb Q,S_0},M_0)=0$, and
$H^1(G_{\mathbb Q,S_0},M_0)$ is an extension of $\mathbb Z_5^2$ by
$\mathbb F_5^3$. Reduction therefore gives the entire five-dimensional
$H^1(G_{\mathbb Q,S_0},\overline M)$. The rational reference class-four
section defines a fixed affine torsor of size $5^5$ for all finite
class-four extensions of the fixed lower class-two section that are
unramified outside $S_0$.

Every actual integral class-four lift for an E.215 tail reduces into
this same torsor. Its class-three predecessor is automatically the
actual local point class at every support prime, because the local
weight-three first cohomology is zero. At a fresh prime the actual point
inertia has exponent divisible by five and vanishes in the exponent-five
class-four quotient. After aligning the integral class-three point image,
the remaining class-four difference is central and unramified. At
zero-support places integral tame inertia is zero as well. Thus these
finite reductions have no ramification outside $S_0$.

Conversely, E.215 gives for the actual two-layer module of the changed
integral pair an exact sequence
\[
 0\longrightarrow\mathbb F_5^3\longrightarrow H^1(G_{\mathbb Q},M)
 \longrightarrow\mathbb Z_5^2\longrightarrow0.
\]
Its reduction modulo five injects into finite cohomology. By the preceding
inertia argument it lands in the same fixed $S_0$-space, and its dimension
is five. It therefore fills that entire five-dimensional space. Thus
\emph{every} one of the $5^5$ fixed finite candidates is occupied by an
integral class-four lift for every changed E.215 tail; the image is not an
unknown proper affine subset. This dimension argument does not assume
that the absolute global second cohomology of the actual module vanishes.

At the place two, global central $\mathbb F_5^3$ localizes isomorphically,
whereas local weight-three first cohomology is zero. Requiring the actual
point class at two therefore leaves exactly twenty-five candidates.
At each support place the lower finite class-three geometric stabilizer
is zero, because $p^i-1$ is a five-adic unit for $i=1,2,3$.
The integral local two-layer first cohomology is $\mathbb F_5^3$ and
injects into its finite reduction. Hence a pinned finite candidate passes
its local finite point test if and only if the corresponding integral
class-four lift passes the actual local point test. The same twenty-five
candidates, with their actual lower frames retained, decide class-four
support matching for every E.215 tail.

The pinned catalog maps bijectively to the two odd-weight-three
parameters modulo five. A coefficient deformation that retains the
literal class-three predecessor modulo five therefore retains one
particular catalog member. Higher coefficient digits cannot repair a
mismatch already visible for that member at an old support prime.
Changing the initial integral odd-layer twist may select another member;
the simultaneous support conditions are still a separate requirement.

\subsubsection*{A nonempty class-four-matched subfamily: the actual gluing}
We prove that the rational-reference member of this catalog can be
retained throughout a selected E.215 recursion. This is a finite-depth
existence assertion and does not constitute progress on the all-depth
question under the task's stated criterion.

At a particular arithmetic step let $F$ be the full finite cyclotomic
CRT field, enlarged to include the needed roots of unity and scalar
arithmetic marks. Let $L_2,L_3$ be the already retained actual two- and
three-primary point-section fields. Let $L_5$ contain the rational-base
class-four exponent-five point section, its actual arithmetic action,
and its finite cusp-path data. Over $F$ the groups of $L_2F,L_3F,L_5F$
are respectively two-, three-, and five-groups. For $L_5$, the action
on the mod-five Frattini quotient is scalar with prime-to-five image
$C_4$, absorbed by $F$, and the kernel of the automorphism map of a
finite five-group to its Frattini automorphisms is a five-group.
The finite cusp-path data adds a five-group. The other coefficients
satisfy the analogous assertion. Consequently the three extensions are
linearly disjoint over $F$.

The actual arithmetic five-data is unramified outside five. The pro-five
tame inertia at 83 and 173 fixes this data and fixes $F$: their local
cyclotomic ramification indices have no factor five, since $82$ and
$172$ are prime to five. Under the rational point section their
abelian vectors are $(1,0)$ and $(0,1)$. The Frattini criterion therefore
gives the full finite geometric five-group over the arithmetic data
and $F$, and hence also over the retained two-/three-primary fields.
Above any compatible arithmetic Frobenius with the fixed CRT restriction
to $F$, there is consequently an actual point-field Frobenius having the
desired cusp-zero, cusp-one, or cusp-infinity graph. The arithmetic
five-component is chosen above the CRT restriction; it is not independently
prescribed in a possibly incompatible compositum.

Let $R=F(w^{1/h})$ be the complementary radical required at the new
prime. In the three directions $w$ is respectively $c/b,c/a,-b/a$.
One of its private valuations at 83 or173 is $1$ or $-1$. The five-part
of $\operatorname{Gal}(R/F)$ is therefore the full cyclic group of order
$5^{v_5(h)}$: a cyclotomic extension has ramification index prime to five
at that private prime, so $w$ cannot become a fifth power in $F$.
Since $R/F$ is cyclic, its intersection
\[
                  C=R\cap(L_2L_3L_5F)
\]
decomposes into primary factors $C_2,C_3,C_5$, with $C_i\subset L_iF$.
The original E.172/E.170/E.128 prescription already glues the two- and
three-primary factors and every previous radical condition.

At the prescribed new prime $q$, the multiplier $\chi(q)$ is primitive
modulo five. Thus $\chi(q)-1$ is invertible on every cyclic five-primary
radical torsor. The prescribed $L_5$ Frobenius has a unique fixed root on
the $C_5$ torsor. Start from the previously compatible fixed-root action
on $R$ and conjugate it by a five-primary radical translation to make
this fixed root agree. Multiplication by $\chi(q)-1$ is invertible,
so the required translation exists. Its prime-to-five restrictions are
unchanged. Two affine actions with the same multiplier and the same fixed
root agree, which makes the two actions identical on $C_5$ without
altering their already equal actions on $C_2,C_3$. The compatible actions
therefore glue to an automorphism of the actual joint field.

We spell out the lower marking, since uniqueness alone would not identify
two differently named roots. The retained literal abelian Kummer equality
identifies the current complementary fifth-root mark: current coordinates
differ from the rational reference by rational fifth powers. The chosen
exact cusp graph has zero five-primary abelian Frobenius coordinates and
therefore fixes this named root. At infinity $-1$ has the compatible
rational fifth root $-1$. First arrange, by a five-primary translation,
that the old compatible radical Frobenius fixes the same named root;
this changes no $F$ or $L_2,L_3$ data. The two affine offsets are now zero
in that lower quotient. Their remaining difference lies in its kernel,
and division by the unit $\chi(q)-1$ keeps the final correcting
translation in that kernel. Thus the literal lower mark is preserved.
If only existence of a radical root was prescribed, no named-root
condition is necessary.

The adjustment changes no rational CRT residue, no chosen old local
unit root, no minimal shell modulo twenty-five, and no two-/three-primary
Frobenius class. The new-prime condition is the existence of an $h$-th
root of $w$; relabeling its five-primary coordinate does not change this
condition. Every prime-to-five radical coordinate remains fixed, including
the negative infinity radical's eighth-root mark from E.170. Thus every
literal E.215 ordinary product-factorization condition is preserved.

For a later active coefficient $\ell\ge7$, include this same five-field
and the prime-to-$\ell$ complementary radical in E.173's retained
arithmetic field. Its commutator has order prime to $\ell$: the old
two-/three-/five-primary factors have only those prime divisors, the
radical factor over $F$ has degree prime to $\ell$, and $F/\mathbb Q$
is abelian. This preserves the active geometric image. Crucially, the
preceding fixed-root construction supplies the actual compatible
retained Frobenius \emph{before} E.153 is invoked; it is not inferred
from group orders. E.153's remaining first-coordinate reciprocity
hypothesis is unchanged, because the prescribed old first coordinates
still come from the single current rational $A$-prefix. The active
$\ell$-primary unit-root condition is automatic since $q-1$ is an
$\ell$-adic unit. This verifies the insertion in the ordinary advance
and every once-only high-coefficient repair.

Run E.215 with these enlarged finite data. Its convergence, frozen
integer valuations, ordinary local identities, all-prime class-two
matches, and single $\mathbb Z/24$ normalization remain valid. At the
private places two, 83, 173, the rational-reference finite class-four
candidate is still the actual point class: valuations are unchanged
and the local unit groups have trivial pro-five completion. At a fresh
support prime the exponent is divisible by five, so point inertia
vanishes in the exponent-five class-four quotient, and the prescribed
cusp Frobenius is precisely its finite point class. The reference
candidate therefore passes every finite support test. The catalog
surjectivity and local injection proved above give an integral global
class-four lift realizing it and matching every off-diagonal support
point. No point condition at five, class-eight matching, or all-depth
tower is asserted.

\subsubsection*{Why this reference argument does not iterate to the whole tower}
Fix one fresh support prime $q$ and let its nonzero frozen ordinary
valuation in a coordinate be $m$. Its five-adic valuation $a=v_5(m)$
is finite. The rational reference coordinate has valuation zero at $q$.
The valuation homomorphism to $\mathbb Z/5^{a+1}$ therefore distinguishes
the two Kummer classes: their difference is $m\not\equiv0$ modulo
$5^{a+1}$. Thus the literal global abelian agreement needed to identify
all the finite twisted modules with one rational-reference module holds
only to bounded coefficient precision. It cannot be assumed at every
later depth. Replacing the reference by a rational arithmetic prefix
does not remedy this: its two rational coordinates generally do not
sum to one, so the construction supplies no rational-point section for
that prefix. In addition E.218 shows that retaining the full interior
reference Frobenius through all pro-five quotients is impossible at
this fixed boundary-support prime. These are limits of the reference
method, not nonexistence statements for arbitrary higher choices.
